% -----------------------------
% 執筆にあたっての留意点
% -----------------------------
%
% 実験・調査の章は，提案事項や仮説を立証するための評価実験，調査の方法を述べる章です．名前のつけ方は内容や分野によって異なります．仮説の立証方法が実験であれば「実験」「ユーザ実験」など，アンケートなどの調査であれば「調査」などが章の名前として考えられます．理学などの分野では，単に「方法」とする場合もあります．
%
% 「実験」の章では，前章までに設定した仮説を検証するための流れを説明してください．一般的には以下の項目に関して，節を立てて説明するとよいでしょう：
%
% * 実験協力者（実験協力者の属性，集め方）
% * 実験手順（時系列に沿って説明）
% * 実験装置・環境（どのような装置を用いたのか，どのような環境で実験をおこなったか，比較手法は何か）
% * 評価・測定指標（仮説と対応づける）
%
% 重要なことは，論文の読者が同じ実験を再現できるよう，必要な情報を正確かつ十分に書くことです．再現できない実験・調査を書いている場合は，反証可能性がないので「トンデモ研究」と見なされます．
%
%
% -----------------------------
% 補足（参考：how-to-xx/how-to-write-a-paper/1-学術論文の書き方.md 3.5節，
%              how-to-xx/how-to-write-a-paper/3-国際会議論文の傾向と対策.md 1章）
% -----------------------------
%
% ■ 検証したいことを，RQ（リサーチクエスチョン）や仮説として明示してから実験を設計します
% 「RQ1: 〜か？」「仮説1: 〜だろう」と番号を振っておくと，結果・考察の章と対応づけやすくなります．
% 測定指標は，必ずどれかのRQ・仮説に紐づけてください．どのRQにも紐づかない指標は，測る必要がありません．
%
% ■ 実験設計では「何を測るか」と「何を示すか」を分けて考えます
% * 何を測るか：実行時間，重要なパラメータへの感度，スケーラビリティ（データ量や問題の複雑さ），主観評価 など
% * 何を示すか：絶対性能（実用に足るか），素朴な手法との比較，既存手法との比較，
%               提案手法の構成要素ごとの寄与（アブレーション）
%
% ■ 自分の手法を良く見せるためだけの実験をしてはいけません
% 有利な条件だけを選んだ実験や，弱すぎるベースラインとの比較は，査読者にすぐ見抜かれます．
% 比較対象は「現時点で最も強い相手」を選んでください．また，なぜそのベースラインを選んだのかの
% 妥当性は，関連研究の章で説明できるようにしておきます．
%
% ■ 人を対象とする実験・調査では，倫理的な手続きを必ず記述します
% 倫理審査（所属機関の研究倫理委員会など）の承認，実験協力者への説明と同意の取り方，謝礼の有無と金額，
% 取得したデータの匿名化と管理方法を，この章に書いてください．国際会議では必須の記述です．
%
% ■ 分析方法も「再現できる粒度」で書きます
% * 定量分析：用いる検定の名前と，その前提（正規性，等分散性など）を確認したかどうか
% * 定性分析：thematic analysisなどの手法名と手順．複数人でコーディングした場合はその一致度，
%             1人で行った場合はその妥当性の説明
%
% ■ 実験の実施前に，この章を書き上げてしまうことを勧めます
% 実験の手順・条件・測定指標を文章にすると，設計の穴（統制していない要因，仮説と対応しない指標など）が
% 実施前に見つかります．実験をやり直す手間を考えれば，先に書いておく方がはるかに安上がりです．

% ----------------------
% 実験・調査
% ----------------------
\section{実験}